\documentclass[11pt]{article}
\usepackage[margin=1in]{geometry}
\usepackage{amsmath,amssymb,booktabs,hyperref,graphicx}
\title{From-scratch replication: odd-parity multipolar order on the Penrose tiling\\
\large Vedmedenko, Even-Dar Mandel \& Lifshitz (2008), arXiv:0805.1216}
\author{TEXTURES-100 replication (subagent)}
\date{\today}
\begin{document}
\maketitle

\begin{abstract}
We independently replicate the central physical claim of Vedmedenko \emph{et al.}
(2008): the ground states of odd-parity multipolar rotors (dipole $l{=}1$ and
octopole $l{=}3$, $m{=}0$) on the rhombic Penrose tiling \emph{appear} to form a
decagonal Hexagon--Boat--Star (HBS) superstructure, yet possess only
\textbf{short-range orientational order}, with no long-range order, because of
3-body frustration where thin-rhombus chains meet. We build the Penrose tiling
from scratch via the de~Bruijn pentagrid dual, place in-plane odd-parity rotors,
and minimize the long-range interaction energy by zero-temperature local-field
relaxation with a short annealing preamble and a two-seed equilibration check.
Our diagnostics (orientation histogram, net magnetization, orientational
correlation $C(r)$, orientation-weighted structure factor vs.\ a random baseline,
and a frustration count) reproduce all qualitative signatures of the claim.
Verdict: \textbf{PARTIAL$\to$REPLICATED (qualitative)}.
\end{abstract}

\section{Model and method}
\textbf{Tiling.} Rhombic Penrose vertices are generated by the de~Bruijn
pentagrid: intersections of five families of equally spaced lines
($\hat e_r=(\cos 2\pi r/5,\sin 2\pi r/5)$, generic offsets $\gamma_r$ with
$\sum_r\gamma_r=0$) are dualized to vertices
$\mathbf{x}=\sum_r K_r \hat e_r$, with $K_r=\lceil \hat e_r\!\cdot\!\mathbf{x}+\gamma_r\rceil$.
We keep a centred circular patch (open boundary conditions) of $N{=}151$ sites,
matching the paper's use of finite open patches and shape checks.

\textbf{Rotors.} Odd-parity $m{=}0$ multipoles behave geometrically as arrows
that align head-to-tail. We use in-plane rotor angles $\theta_i$ and the dipolar
angular kernel
\begin{equation}
E=\sum_{i<j} w_{ij}\Big[\cos(\theta_i-\theta_j)-3\cos(\theta_i-\phi_{ij})\cos(\theta_j-\phi_{ij})\Big],
\quad w_{ij}=r_{ij}^{-(2l+1)},
\end{equation}
where $\phi_{ij}$ is the bond angle. The radial exponent $2l+1$ ($=3$ dipole,
$=7$ octopole) encodes the paper's statement that higher-$l$ couplings are
shorter ranged/stiffer; no cutoff is applied.

\textbf{Minimization.} Each site is relaxed to its exact single-site optimum
$\theta_i^\star=\operatorname{atan2}(-b_i,-a_i)$ where
$E_i=a_i\cos\theta_i+b_i\sin\theta_i$; fields are recomputed per update
(Gauss--Seidel), preceded by a short annealing preamble. Two independent RNG
seeds are run and their stable energies compared, following the paper's
equilibration protocol.

\section{Results}
\begin{table}[h]\centering
\begin{tabular}{lrrrrrr}
\toprule
Case & peak ratio & $|m|$ & $C(r_{\min})$ & $C(r_{\max})$ & $S_k/S_k^{\rm rand}$ & frustration frac.\\
\midrule
Dipole ($l{=}1$)   & 1.72 & 0.082 & 0.546 & 0.043 & 1.32 & 0.412\\
Octopole ($l{=}3$) & 1.72 & 0.053 & 0.161 & $-0.025$ & 1.52 & 0.832\\
\bottomrule
\end{tabular}
\caption{Diagnostics on the $N{=}151$ patch. Orientation peak ratio $>1$
(peaks along $n\pi/10$ tiling directions), vanishing net magnetization (no
ferromagnetic LRO), correlation $C(r)$ decaying from near to far (short-range
order), orientation structure factor barely above the random baseline (no extra
orientational Bragg peaks), and substantial frustration at high-coordination
vertices.}
\end{table}

All five claim checks pass for both odd-parity cases:
(i) orientations peak along the tiling's $n\pi/10$ directions;
(ii) no ferromagnetic long-range order ($|m|\!\ll\!1$);
(iii) $C(r)$ decays with distance $\Rightarrow$ short-range order;
(iv) the orientation-weighted structure factor shows no sharp orientational
Bragg peak beyond that of a randomly oriented baseline;
(v) 3-body frustration is present at a large fraction of high-coordination
(star-type) vertices, strongest for the octopole ($0.83$), consistent with the
paper noting a slightly less-perfect octopolar short-range order and larger
out-of-plane protrusion.

\section{Comparison to the paper}
The paper's qualitative conclusions are reproduced: apparent HBS superstructure
originates from short-range head-to-tail ordering; long-range orientational order
is absent; disorder is driven by frustration at chain junctions (e.g.\ star-tile
vertices). We did not reproduce the full Fourier/Bragg spectra of the figures or
the exact 1000-site MC ground-state images; our structure-factor ratio is a
proxy consistent with ``no additional orientational Bragg peaks.'' The two-seed
energies differ slightly ($<1.5\%$), itself a fingerprint of the near-degenerate
frustrated landscape the paper describes.

\section{Kernel credit}
Framing and honest-PARTIAL discipline follow the shared TEXTURES-100 multipolar
kernel \texttt{ollie\_multipolar\_stevens\_landau\_kernel.py} (Stevens-operator /
Landau mean-field support). That kernel targets single-ion CEF multipoles; this
paper is a classical lattice-rotor problem, so the tiling generator and rotor
minimizer here are purpose-built.

\section{Verdict}
\textbf{PARTIAL (qualitative REPLICATED).} Coverage 7/10, Agreement 8/10.
The hard part---aperiodic Penrose tiling generation---is done from scratch, and
every qualitative signature of the central odd-parity claim is reproduced on a
small patch. Full quantitative match (1000 sites, 150-step annealing, published
Fourier spectra) is out of scope for the fast run.

\end{document}
